% 2020 北京卷 第 15 题：甲、乙两企业污水排放量 W 与时间 t 的关系曲线
% 来源 PDF：普通高考/2020/2020北京.pdf 第 1 页（A3 横排），印面 1113，第 15 题题干下方
% 重绘依据：原页该图为矢量路径，已由 pdftocairo -svg 解出四条折线（两企业各被 t2 处的
%   交点切成左右两段），按坐标轴端点（O、t 轴右端、W 轴上端）线性换算为本图坐标，
%   采样后以 plot 还原，形状与原页一致。
% 原页要素：横轴 t、纵轴 W；水平虚线「污水达标排放量」；t 轴上 t1、t2、t3 三条竖直虚线
%   （t1 在甲曲线开始明显下降之后、t2 在两曲线交点、t3 在两曲线都降到达标线以下之后）；
%   曲线起点处标「甲企业」（上）、「乙企业」（下），交点右下方标「乙企业」（上）、
%   「甲企业」（下）。
\begin{tikzpicture}[line width=0.7pt, font=\small,
  x=0.50cm, y=0.50cm, line cap=round, line join=round]

  % 坐标轴（W 轴比原页留高一点，供曲线标注）
  \draw[->] (0,0) -- (15.9,0) node[below, inner sep=5pt] {$t$};
  \draw[->] (0,0) -- (0,13.4) node[right, inner sep=3pt] {$W$};
  \node[below left, inner sep=2pt] at (0,0) {$O$};

  % 污水达标排放量（水平虚线；标注按原页置于纵轴左侧）
  \draw[dashed] (0,2.92) -- (15.1,2.92);
  \node[left, inner sep=2pt] at (-0.15,2.92) {污水达标排放量};

  % 甲企业曲线（左段：起点最高，先略升后下降；右段：过交点后降得最快）
  \draw plot[smooth] coordinates {
    (0.00,8.76) (0.82,8.86) (1.63,8.87) (2.45,8.79) (3.26,8.62)
    (4.07,8.37) (4.85,8.02) (5.62,7.60) (6.37,7.08) (7.10,6.49)
    (7.79,5.82) (8.45,5.08) (9.07,4.26)
  };
  \draw plot[smooth] coordinates {
    (9.07,4.26) (9.45,3.85) (9.83,3.47) (10.23,3.12) (10.64,2.80)
    (11.06,2.51) (11.49,2.26) (11.92,2.04) (12.37,1.85) (12.81,1.70)
    (13.27,1.58) (13.72,1.50) (14.18,1.46)
  };

  % 乙企业曲线（左段：起点较低，单调缓降；右段：过交点后仍在甲之上）
  \draw plot[smooth] coordinates {
    (0.00,6.60) (0.77,6.59) (1.54,6.55) (2.31,6.47) (3.08,6.36)
    (3.84,6.22) (4.60,6.04) (5.36,5.83) (6.11,5.58) (6.86,5.30)
    (7.61,4.99) (8.34,4.64) (9.07,4.26)
  };
  \draw plot[smooth] coordinates {
    (9.07,4.26) (9.47,3.98) (9.88,3.71) (10.29,3.47) (10.71,3.24)
    (11.13,3.05) (11.56,2.87) (11.99,2.72) (12.43,2.60) (12.86,2.50)
    (13.30,2.42) (13.74,2.37) (14.18,2.34)
  };

  % t1、t2、t3 三条竖直虚线（高度分别止于对应曲线）
  \draw[dashed] (3.89,0) -- (3.89,8.58);
  \draw[dashed] (9.07,0) -- (9.07,4.26);
  \draw[dashed] (13.80,0) -- (13.80,2.39);
  \node[below, inner sep=4pt] at (3.89,0) {$t_1$};
  \node[below, inner sep=4pt] at (9.07,0) {$t_2$};
  \node[below, inner sep=4pt] at (13.80,0) {$t_3$};

  % 曲线标注（按原页位置：起点标注在曲线上方，交点后的标注在右侧）
  \node[right, inner sep=2pt] at (0.45,10.18) {甲企业};
  \node[right, inner sep=2pt] at (0.45,7.85) {乙企业};
  \node[right, inner sep=2pt] at (10.95,4.62) {乙企业};
  \node[right, inner sep=2pt] at (10.95,1.52) {甲企业};
\end{tikzpicture}
